\documentclass[11pt]{article}
\usepackage[T1]{fontenc}
\usepackage{lmodern,amsmath,amssymb,amsthm,booktabs,array}
\usepackage[margin=1in]{geometry}
\usepackage[colorlinks=true,allcolors=blue]{hyperref}
\newtheorem{theorem}{Theorem}
\newtheorem{corollary}[theorem]{Corollary}
\newtheorem{proposition}[theorem]{Proposition}
\newcommand{\F}{\mathbb F_2}
\newcommand{\HH}{H}
\title{A counterexample to the Shannon conjecture following Eq.~(69)\newline
in Tarabunga et al., Section V.C}
\author{Haobo Ma\\{\small ChronoAI Pte Ltd, Singapore}
\and Wenlin Zhang\\{\small National University of Singapore}\\{\small ChronoAI Pte Ltd, Singapore}}
\date{November 2, 2026}
\begin{document}
\maketitle
\begin{abstract}
The Shannon conjecture following Eq.~(69) in Section V.C of
Tarabunga et al., arXiv:\allowbreak2411.05766v1, is false. An explicit three-qubit
state has stabilizer-basis entropy $\log 4$, whereas the proposed
absolute-amplitude lower bound is $\log16-(9/8)\log3>\log4$.
The same construction refutes the bound for nonnegative amplitudes,
and continuity gives strictly positive examples. It extends to every
$N\ge3$ and to a range of R\'enyi indices containing one.
The collision entropy of the comparison distribution remains a valid
Shannon lower bound for arbitrary affine supports. For amplitudes of constant
modulus on each affine coset, the original bound holds at every R\'enyi index;
this includes the source's Figure~1 product states.
\end{abstract}

\section{The conjecture and its specialization}
Tarabunga, Frau, Haug, Tirrito and Piroli conjecture the Shannon extension
of their branch-and-bound inequality after Eq.~(69),
Section V.C of arXiv:2411.05766v1~\cite{posing}: the equation is on p.~9
and the explicit Shannon conjecture is on p.~10.
For a fixed binary support matrix $R$, it asserts
\begin{equation}\label{eq:conjecture}
 \min_{Q',c}\HH(p_{Q',c,R})\ \ge\ \HH(q_R).
\end{equation}
Here $Q'$ is strictly upper triangular, $c$ is binary, and
$q_R$ is obtained by replacing the computational amplitudes of the state
by their absolute values and setting $Q'=c=0$.
The source proves the corresponding inequality for integer
R\'enyi index greater than one. The journal version~\cite{journal}
retains the conjecture in Section~5.3, after Eq.~(74).
Unless stated otherwise, equation and appendix locators here refer to
the specified arXiv version.
We use natural logarithms and $0\log0=0$.

It suffices to take full support, $k=N$ and $R=I_N$.
For $d,x\in\F^N$ define the orthonormal stabilizer basis
\begin{equation}\label{eq:basis}
 b_{Q',c,d}(x)=2^{-N/2}(-1)^{d\cdot x+x^TQ'x}i^{\sum_j c_jx_j}.
\end{equation}
The dot product and quadratic exponent are binary; the exponent of $i$
is an ordinary integer sum. Orthogonality follows by summing the
characters $(-1)^{(d+e)\cdot x}$; the other factors have modulus one.
This is the integer-sum convention of~\cite{lean}. Interpreting the
source's $c^Tx$ as a binary-field dot product instead gives a parity
exponent. Section~\ref{sec:real} gives the quadratic relabeling between
these parametrizations; the $c=0$ counterexample is identical in both.
For a normalized amplitude vector $\psi$, set
$p_{Q',c}(d)=|\sum_x\overline{b_{Q',c,d}(x)}\psi(x)|^2$
and $q(d)=|\sum_xb_{0,0,d}(x)|\psi(x)||^2$.
Thus $q$ and $p_{0,0}$ are squared Walsh transforms of $|\psi|$ and $\psi$.

\section{An exact counterexample}
The signed witness below is encoded in the pinned Lean source~\cite{lean},
whose declaration \texttt{result} refutes the full-support Shannon bound.
We give its entropy calculation and consequences explicitly.
\begin{theorem}\label{thm:main}
In the computational order $000,001,\ldots,111$, let
\[
 \psi=\frac1{\sqrt8}(1,0,2,-1,1,0,0,1).
\]
For $R=I_3$ and $Q'=c=0$, the comparison and overlap distributions satisfy
\[
 q=\frac1{16}(9,1,1,1,1,1,1,1),\qquad
 p=\frac14(1,1,0,0,0,1,0,1).
\]
Consequently \eqref{eq:conjecture} is false.
\end{theorem}
\begin{proof}
The squared coefficients sum to eight. If $W_{d,x}=(-1)^{d\cdot x}$,
direct integer multiplication gives
\begin{align*}
 W(1,0,2,-1,1,0,0,1)^T&=(4,4,0,0,0,4,0,-4)^T,\\
 W(1,0,2,1,1,0,0,1)^T&=(6,2,-2,2,2,2,-2,-2)^T.
\end{align*}
The overlaps are these entries divided by eight, proving the formulas.
The entropies are
\[
 \HH(p)=\log4,\qquad
 \HH(q)=\log16-\frac98\log3.
\]
Their difference is
\begin{equation}\label{eq:gap}
 \HH(q)-\HH(p)=\frac18\log\frac{2^{16}}{3^9}>0,
 \qquad 65536>19683.
\end{equation}
Since the minimum in \eqref{eq:conjecture} is at most $\HH(p)$,
one admissible basis suffices for the contradiction.
\end{proof}
The gap is approximately $0.1503555364$ nats, or $0.2169171867$ bits.
The example establishes failure of the proposed lower bound;
it does not assert a value for the minimization over all support matrices $R$.

\begin{corollary}\label{cor:positive}
The Shannon lower bound fails even for nonnegative amplitudes,
and also for some strictly positive amplitudes.
\end{corollary}
\begin{proof}
Put $a=|\psi|$ and take $c=0$ with quadratic phase
$x^TQ'x=x_0x_1+x_1x_2$. On the support of $a$, its sign is negative
at $011$ and positive at the other four occupied strings.
Thus $(-1)^{x^TQ'x}a(x)=\psi(x)$, so this basis has exactly the
distribution $p$ of Theorem~\ref{thm:main}, while its comparison is $q$.
For strict positivity, replace the three zero coefficients of $a$ by
$\varepsilon>0$ and normalize. Both distributions depend continuously
on $\varepsilon$, and so does entropy because $-u\log u$ extends
continuously to $u=0$. The strictly positive gap at $\varepsilon=0$
therefore persists for all sufficiently small positive $\varepsilon$.
\end{proof}

An explicit strictly positive instance is
\begin{equation}\label{eq:strict}
 a_+=\frac1{\sqrt{80003}}(100,1,200,100,100,1,1,100).
\end{equation}
With the same quadratic phase, the integer Walsh vectors for the comparison
and overlap distributions are, respectively,
\[
 (603,199,-199,197,199,199,-199,-199),\quad
 (401,397,3,-1,1,401,-1,-401).
\]
Squaring each entry and dividing by $8\cdot80003$ specifies both
distributions exactly. The finite check in \texttt{verify.py} certifies
$137/1000<\HH(q)-\HH(p)<139/1000$ by rational logarithm bounds.
For $1\le r\le2$ it uses
\[
 \log r=2\sum_{j=0}^{L-1}\frac{t^{2j+1}}{2j+1}+E_L,
 \quad t=\frac{r-1}{r+1},\quad
 0\le E_L\le\frac{2t^{2L+1}}{(2L+1)(1-t^2)},
\]
with $L=70$ and reduction by powers of two.

\begin{corollary}\label{cor:extensions}
Counterexamples exist for every $N\ge3$. For the distributions in
Theorem~\ref{thm:main}, $S_\alpha(q)>S_\alpha(p)$ holds exactly for
$0\le\alpha<\alpha_*$, where $\alpha_*>1$ is the unique root above one of
\begin{equation}\label{eq:root}
 (9/4)^{\alpha_*}+7\,4^{-\alpha_*}=4.
\end{equation}
Here $S_0$ is log support size and $S_1=\HH$;
numerically $\alpha_*\simeq1.262623$.
No dimension-independent additive correction repairs the Shannon bound.
\end{corollary}
\begin{proof}
Tensor $\psi$ with $|0\rangle^{\otimes(N-3)}$ and use the product Walsh
basis. Both outcome distributions acquire a uniform factor of size
$2^{N-3}$, adding $(N-3)\log2$ to each R\'enyi entropy, including
indices zero and one. The gap is unchanged.
Strictly positive examples in every such dimension follow by tensoring
$a_+$ from \eqref{eq:strict} with $|+\rangle^{\otimes(N-3)}$ instead:
the extra computational amplitudes are positive and the extra Walsh
outcome is a point mass, so the entropy gap is unchanged.

For $\alpha>0$, $\alpha\ne1$, the difference of power sums has the sign of
$g(\alpha)=(9/4)^\alpha+7\,4^{-\alpha}-4$.
This function is strictly convex, $g(0)=4$, $g(1)=0$,
and $g(2)=3/2>0$. Moreover
$g'(1)=(9\log3-16\log2)/2<0$ by \eqref{eq:gap}.
Strict convexity gives exactly one further root $\alpha_*\in(1,2)$,
with $g>0$ below one, $g<0$ between one and $\alpha_*$,
and $g>0$ above $\alpha_*$. Division by $1-\alpha$ in
$S_\alpha(r)=\log(\sum_dr(d)^\alpha)/(1-\alpha)$ gives the stated signs.
At zero the support sizes are eight and four; at one use Theorem~\ref{thm:main}.
Finally, tensoring $b$ copies gives $p^{\otimes b}$ and $q^{\otimes b}$
in an admissible product basis, with Shannon gap $b\delta$, where
$\delta$ is \eqref{eq:gap}. Thus no universal constant, or even an
$o(N)$ term, can bound the defect in dimension $N=3b$.
The same amplification applies to the strictly positive example.
\end{proof}

\section{The real-amplitude reduction}
\label{sec:real}
Appendix C.2 of~\cite{posing} reduces global minimization to $c=0$
by pairing complex overlaps and comparing them with real overlaps.
In the integer-sum convention of \eqref{eq:basis}, that argument
requires a quadratic phase conversion. To see why, consider the real
two-qubit state $\eta=(1,1,1,-1)/2$, with $Q'=0$ and $R=I_2$.
At $c=(0,0)$, the outcome distribution is uniform on four points;
at $c=(1,1)$, it is $(1/2,0,0,1/2)$.
Their R\'enyi entropies are $\log4$ and $\log2$ for every $\alpha\ge0$.
Thus the integer-sum fixed-$Q'$ analogue of (C3) is false.
This is a parametrization distinction, not a counterexample to the
binary-parity form of (C3). The global reduction to real bases has the
following complete proof in the integer-sum convention.

\begin{proposition}\label{prop:real}
For a real amplitude vector and fixed $R$, the minimum over $(Q',c)$
is attained at $c=0$, for every $0\le\alpha\le\infty$.
More precisely, define the strictly upper-triangular matrix
$(B_c)_{ij}=c_ic_j$ for $i<j$. Then
\[
 S_\alpha(p_{Q'+B_c,0,R})\le S_\alpha(p_{Q',c,R}).
\]
\end{proposition}
\begin{proof}
For binary $x$, let $s=\sum_jc_jx_j$ as an integer.
Since $\lfloor s/2\rfloor\equiv\binom{s}{2}\pmod2$,
\[
 (-i)^s=(-1)^{x^TB_cx}(-i)^{c\cdot x},
\]
where the last exponent $c\cdot x$ is reduced to zero or one.
If $c=0$, the claim is equality. Otherwise pair each diagonal label $d$
with $d+c$, separately in each coset $t$. After absorbing
$(-1)^{x^TB_cx}$ into $Q'$, split the real Walsh sum into its even- and
odd-$c\cdot x$ parts, denoted $\beta$ and $\gamma$, including the
factor $2^{-k/2}$. The original complex overlaps are
$\beta-i\gamma$ and $\beta+i\gamma$, giving paired probabilities
\[
 (\beta^2+\gamma^2,\,\beta^2+\gamma^2),
\]
whereas the adjusted real basis has
\[
 ((\beta+\gamma)^2,\,(\beta-\gamma)^2).
\]
The pairs have the same sum. The second pair is at least as imbalanced:
convexity of $u^\alpha$ for $\alpha>1$ gives a larger power sum,
and concavity for $0<\alpha<1$ gives a smaller one.
For $\alpha=1$, concavity of $-u\log u$ gives smaller entropy;
for $\alpha=0$, the support cannot increase. Summing over all pairs
and using the sign of $1-\alpha$ proves the inequality.
At $\alpha=\infty$, the largest entry of each real pair is at least
as large as either entry of the equal complex pair.
Each adjusted $Q'+B_c$ is admissible. Since $c=0$ is already an allowed
choice, the two global minima are equal.
\end{proof}
In particular, an imaginary linear phase cannot simply be discarded while
holding the quadratic phase fixed; its parity contributes $B_c$.
If the exponent of $i$ is instead defined as binary parity from the outset,
this quadratic adjustment is already absorbed in the parametrization.

\begin{corollary}\label{cor:minimum}
For $a=|\psi|$ from Theorem~\ref{thm:main}, and every
$0\le\alpha\le\infty$,
\[
 \min_{Q',c} S_\alpha(p^a_{Q',c,I_3})
 =\min\{\log4,S_\alpha(q)\}.
\]
\end{corollary}
\begin{proof}
Proposition~\ref{prop:real} reduces the minimum to $c=0$.
On the occupied strings, each quadratic phase yields coefficients
$(1,0,2,s,1,0,0,t)/\sqrt8$, with $s,t\in\{1,-1\}$;
every sign pair occurs. If $s=t$, the squared Walsh transform is a
permutation of $q$; if $s=-t$, it is uniform on four points.
These four integer transforms determine the displayed minimum.
\end{proof}

\section{A valid Shannon lower bound}
The integer-index argument in Appendix C.1 of~\cite{posing} gives a
replacement that does not require the Shannon conjecture.
\begin{proposition}\label{prop:safe}
For every normalized state, support matrix $R$ of rank $k$, and
admissible $Q',c$, one has
\begin{equation}\label{eq:safe}
 \HH(p_{Q',c,R})\ \ge\ -\log\sum_{d,t}q_R(d,t)^2.
\end{equation}
\end{proposition}
\begin{proof}
Write $M=2^k$, choose one representative $t$ of each coset of
$\operatorname{Im}R$, and put
\[
z_t(x)=(-1)^{x^TQ'x}(-i)^{\sum_jc_jx_j}\psi(Rx+t).
\]
Then $p(d,t)=M^{-1}|\sum_x(-1)^{d\cdot x}z_t(x)|^2$,
while $q_R$ replaces $z_t$ by $|z_t|$.
Character orthogonality, valid also at $k=0$, gives
\begin{align*}
 \sum_{d,t}p(d,t)^2
 &=\frac1M\sum_t\sum_{x_1+x_2=x_3+x_4}
 z_t(x_1)z_t(x_2)\overline{z_t(x_3)z_t(x_4)}\\
 &\le\frac1M\sum_t\sum_{x_1+x_2=x_3+x_4}
 |z_t(x_1)z_t(x_2)z_t(x_3)z_t(x_4)|
 =\sum_{d,t}q_R(d,t)^2.
\end{align*}
The inequality is the triangle inequality; the left-hand sum is real
and nonnegative. Both distributions are normalized, since the cosets
partition the computational basis and each Walsh transform is unitary.
Jensen's inequality for $-\log$ with weights $p(d,t)$ yields
$\HH(p)\ge-\log\sum p(d,t)^2$, proving \eqref{eq:safe}.
\end{proof}
The same expansion for integer $m\ge2$ has factor $M^{1-m}$ and
constraint $x_1+\cdots+x_m=y_1+\cdots+y_m$.
Taking absolute values proves $\sum p^m\le\sum q_R^m$, hence
$S_m(p)\ge S_m(q_R)$, the source's integer-index result.
For Theorem~\ref{thm:main}, the safe bound is $\log(32/11)$,
because $\sum q^2=11/32$. Thus a Shannon branch can be pruned safely
when this collision-entropy bound exceeds the incumbent entropy;
the larger conjectural quantity $\HH(q_R)$ is not a universal certificate.

There is a stronger coset-conditioned version. Let
$w_t=\sum_d p(d,t)=\sum_d q_R(d,t)$ and
$\widehat q_t(d)=q_R(d,t)/w_t$ for $w_t>0$.
The same collision argument within each coset and the Shannon chain rule give
\begin{equation}\label{eq:conditioned}
 \HH(p)\ge \HH(w)+\sum_{t:w_t>0}w_tS_2(\widehat q_t)
 \ge S_2(q_R).
\end{equation}
For the last inequality, put $C_t=\sum_dq_R(d,t)^2$.
The middle expression is $-\sum_tw_t\log(C_t/w_t)$, which is at least
$-\log\sum_tC_t$ by Jensen's inequality. Zero-mass cosets contribute zero.

\begin{proposition}\label{prop:flat}
If $|\psi(Rx+t)|=a_t$ is constant in $x$ on every affine coset, then
$S_\alpha(p_{Q',c,R})\ge S_\alpha(q_R)$ for every
$0\le\alpha\le\infty$ and every $Q',c$.
\end{proposition}
\begin{proof}
Character orthogonality gives $q_R(d,t)=w_t\mathbf1_{\{d=0\}}$,
where $w_t=2^ka_t^2$, while $\sum_dp(d,t)=w_t$.
For $\alpha>1$, splitting a mass gives
$\sum_dp(d,t)^\alpha\le w_t^\alpha$; for $0<\alpha<1$ the
inequality reverses. Summing and dividing the logarithm by $1-\alpha$
proves the assertion. At $\alpha=1$, use
$\HH(p)=\HH(w)+\sum_tw_t\HH(p(\cdot\mid t))\ge\HH(w)=\HH(q_R)$.
At zero the support cannot shrink when masses are split; at infinity
$\max_{d,t}p(d,t)\le\max_tw_t$.
\end{proof}
The normalized Figure~1 states of~\cite{posing},
$2^{-N/2}(|0\rangle+e^{i\theta}|1\rangle)^{\otimes N}$ with
$\theta=\pi/4,\pi/8$, have globally constant amplitude modulus.
For every rank-$k$ support matrix, their comparison entropy is
$(N-k)\log2$. Thus the original pruning bound is valid on these
particular inputs at every index. This conclusion concerns the mathematical
bound, not the numerical implementation or any further transformations.

\section{Coverage and verification scope}
The neighboring claims have distinct scopes:
\begin{center}\small
\begin{tabular}{@{}p{.42\linewidth}p{.52\linewidth}@{}}
\toprule
Source item in arXiv:2411.05766v1 & Conclusion here\\
\midrule
Section V.C, conjecture after (69), $\alpha=1$ & Refuted by Theorem~\ref{thm:main}.\\
Same bound at nearby indices and $\alpha=0$ & Refuted for $0\le\alpha<\alpha_*$ by Corollary~\ref{cor:extensions}.\\
Section V.C, Figure~1 product-state families & The original bound is valid at every index by Proposition~\ref{prop:flat}.\\
Appendix C.2, final positive-amplitude observation & At $\alpha=1$, $Q'=0$ need not minimize: Corollary~\ref{cor:positive}, including strictly positive examples. The integer-index statement is preserved.\\
Section V.C and Appendix C.1, integer $\alpha>1$ & Preserved; the moment proof also gives Proposition~\ref{prop:safe}.\\
Appendix C.2, real-amplitude $c=0$ reduction, Eq.~(C3) & Preserved; Proposition~\ref{prop:real} gives the integer-sum phase conversion $Q'\mapsto Q'+B_c$. Its fixed-$Q'$ integer-sum analogue fails for $\eta$.\\
Sections II.D and V.C, additivity discussions; Section VI, questions on strong monotonicity & Not covered. The counterexample alone decides neither.\\
\bottomrule
\end{tabular}
\end{center}
The pinned \href{https://github.com/the-omega-institute/trureturing/blob/122c987d276498783de7d873378a28b27ced8761/D5/S3/Quantum/Magic/StabilizerAsymmetryShannonBoundRefutation.lean}{Lean source}
defines the full-support distributions in \eqref{eq:basis} and its declaration
\texttt{result} proves the negation of the full-support Shannon bound,
using the signed witness of Theorem~\ref{thm:main}.
Corollaries~\ref{cor:positive}, \ref{cor:extensions} and~\ref{cor:minimum},
Propositions~\ref{prop:real}, \ref{prop:safe} and~\ref{prop:flat}, and
\eqref{eq:conditioned} have written proofs here.
The accompanying finite checks establish the displayed transforms,
the strictly positive instance \eqref{eq:strict}, and, by all 64 choices
of $(Q',c)$, the fixed-$R=I_3$ minimum $\log4$ for $|\psi|$.
These checks do not minimize over other support matrices.

\begin{thebibliography}{3}
\bibitem{posing} P.~S.~Tarabunga, M.~Frau, T.~Haug, E.~Tirrito and L.~Piroli,
\emph{A nonstabilizerness monotone from stabilizerness asymmetry},
\href{https://arxiv.org/abs/2411.05766v1}{arXiv:2411.05766v1} (8 November 2024),
Section V.C, pp.~9--10, Eq.~(69) and the subsequent Shannon conjecture;
Appendix C.1--C.2, pp.~15--16.
\bibitem{journal} P.~S.~Tarabunga, M.~Frau, T.~Haug, E.~Tirrito and L.~Piroli,
\emph{A nonstabilizerness monotone from stabilizerness asymmetry},
Quantum Science and Technology \textbf{10}(4), 045026 (2025),
\href{https://doi.org/10.1088/2058-9565/adfd0d}{doi:10.1088/2058-9565/adfd0d}.
\bibitem{lean} Trureturing contributors,
\emph{StabilizerAsymmetryShannonBoundRefutation},
\href{https://github.com/the-omega-institute/trureturing/blob/122c987d276498783de7d873378a28b27ced8761/D5/S3/Quantum/Magic/StabilizerAsymmetryShannonBoundRefutation.lean}{source at commit \texttt{122c987d2764}};
declaration \texttt{D5.S3.Quantum.Magic.}\\
\texttt{StabilizerAsymmetryShannonBoundRefutation.result}.
\end{thebibliography}
\noindent\small\textcopyright\ 2026 Haobo Ma and Wenlin Zhang.
Licensed under \href{https://creativecommons.org/licenses/by/4.0/}{Creative Commons Attribution 4.0 International (CC BY 4.0)}.
\end{document}
